\documentclass[10pt,a4paper]{amsart}
\usepackage[margin=19mm]{geometry}
\usepackage{amsmath,amssymb,mathtools}
\usepackage[scr=rsfs,cal=euler]{mathalfa}
\usepackage{xfrac}
\usepackage[unicode,hidelinks,bookmarks=false]{hyperref}
\newcommand{\ie}{i.e.\ }
\newcommand{\cf}{cf.\ }
\newcommand{\resp}{respectively\ }
\newcommand{\Subsection}{\subsection}
\providecommand{\nobreakdash}{\nobreak\hbox{-}}
\setlength{\emergencystretch}{2em}
\multlinegap0pt
\usepackage{bm}
\usepackage{bbm}
\usepackage{dsfont}
\usepackage{xspace}
\usepackage{cancel}
\usepackage{mathtools}
\usepackage{cleveref}
\usepackage{amsthm}
\makeatletter
\@ifundefined{thm}{\newtheorem{thm}{Theorem}}{}
\@ifundefined{corollary}{\newtheorem{corollary}[thm]{Corollary}}{}
\makeatother
\title[Generalized fruit Diophantine equation and super elliptic cu]{Generalized fruit Diophantine equation and super elliptic curves}
\author{Kalyan Banerjee, Kalyan Chakraborty, Ankita Das}
\date{Source: arXiv:2602.01001v2 (CC BY 4.0)}
\begin{document}
\maketitle

\noindent\emph{Source and licence.} Generalized fruit Diophantine equation and super elliptic curves. Version arXiv:2602.01001v2 (CC BY 4.0), distributed under the Creative Commons Attribution 4.0 International licence, as recorded in the arXiv licence metadata for this exact version and shown on the abstract page https://arxiv.org/abs/2602.01001v2. Full rights notice: licenses/arxiv-2602.01001-CC-BY-4.0.txt.\par\smallskip

\noindent\textit{Editorial scope.} Counted unit: the Corollary whose source label is \texttt{cor:1} in the article (source0.tex lines 563--569, source label \texttt{cor:1}), delivered together with the article's own complete original proof of it (source0.tex lines 571--609, one \texttt{proof} environment of 39 lines). No mathematics is written, repaired or re-derived: statement and proof are the article's own text line for line. Required context retained uncounted: th1 (lines 407--415); eq2 (lines 428--430). Every definition, display and label of those lines is retained unchanged. Omitted from this package and not redistributed: 1--406, 416--427, 431--562, 570--570, 610--1555. Nothing inside a retained excerpt is omitted or altered: every retained body is delivered exactly as the article's lines are written, so no omission receipt of any kind accompanies this package. Citations: the retained excerpts carry no citation command at all, so no bibliography entry of the article is transcribed and no reference is redistributed. Reference adaptation: none was needed and none was made; every reference inside the delivered unit resolves inside it, so no unresolved reference or citation is produced. Printed slips kept literal: no printed word, symbol, hypothesis or number of the counted statement or of its proof was altered. Layout and class: the article's own class is \texttt{amsart} with packages including fullpage, enumerate, setspace, listings, bm, amsmath, amsfonts, amssymb, graphics, amsthm, comment, hyperref, dsfont, mdframed; the wrapper is the lane's pinned \texttt{amsart} editorial wrapper at 10pt on A4, which restates the theorem-like environments the retained text needs and adds the article's own macro definitions that the retained text needs (none), with extra wrapper packages (bm, bbm, dsfont, xspace, cancel, mathtools, cleveref). The delivered numbering is the unit's own. Assets: no retained span contains a figure environment, a graphics-inclusion command, a PDF-page inclusion, a table or a TikZ picture; no image file is copied into this package, the package declares no asset dependency and no image file is redistributed with it. Licence: version arXiv:2602.01001v2 of this article is distributed under the Creative Commons Attribution 4.0 International licence, as recorded in the arXiv licence metadata for this exact version and shown on the abstract page https://arxiv.org/abs/2602.01001v2. A full rights notice is delivered at licenses/arxiv-2602.01001-CC-BY-4.0.txt. Characters: every retained excerpt is pure ASCII, so the delivered engine, XeLaTeX with its default Latin Modern text font, reports no missing character.
\begin{thm}
The equation 
\[ax^{d_1}-y^{d_2}-z^2+xyz-b=0\] has no integer solution for fixed a and b such that,
\[a\equiv 1 \pmod {12}\]
\[b=2^{d_1}a-3\]
\[d_1 \geq 3 , 3\mid {d_1}\]
\[d_2 \geq 2, 2\mid {d_2}\]
\[gcd(d_1, d_2)=1.\]\label{th1}
\end{thm}

\begin{equation}
  y^{d_2}+z^2+b=a\alpha^{d_1}+\alpha yz.  \label{eq2} 
\end{equation}

\begin{corollary}\label{cor:1}
    Let us consider the equation:
    \[ax^{d_1}-y^{d_2}-z^2+xyz-b=0\]
    This has no integer solution $(x,y,z)$ with $x$ being even for fixed integers $a$ and $b$ such that
    \[a\equiv 1 \pmod {12}\]
\[b=2^{d_1}a-3^m\] with $m$ being odd positive integer and $d_1, d_2$ as in \Cref{th1}. 
\end{corollary}

\begin{proof}
    Clearly, we can follow exactly the same steps in Case 1 of the proof of \Cref{th1}.\\
    Let us assume, there exists a solution with $x= \alpha$ even, then we rewrite \eqref{eq2} as
    \begin{equation}
        \left( y^{2s-1}-\frac{\alpha z}{2}\right)^2-\left( \frac{\alpha ^2}{4}-y^{2(s-1)}\right)z^2=y^{2(s-1)}[ a\alpha ^{d_1}-b].  \label{eq8}
    \end{equation}
    Now set,
    \begin{equation*}
    Y= y^{2s-1}-\frac{\alpha z}{2}
    \end{equation*}
    \begin{equation*}
        \beta= \frac{\alpha}{2}
    \end{equation*}
    \begin{equation*}
        k=y^{s-1}
    \end{equation*}
    \begin{equation*}
        Z=z.
    \end{equation*}
    Then \eqref{eq8} can be written as,
    \begin{equation}
    Y^2-\left(\beta ^2-k^2\right )Z^2=k^2\left[ a\alpha ^{d_1}-b\right] = k^2\left[a2^{d_1}\beta ^{d_1}-a2^{d_1}+3\right] \label{eq9}
    \end{equation}
    If $\beta$ is even, then reducing \eqref{eq9} modulo 4 gives,
    \begin{equation*}
        Y^2 \equiv 3^m \equiv 3 \pmod{4} 
    \end{equation*} or 
    \begin{equation}
        Y^2+Z^2 \equiv 3^m \equiv 3 \pmod{4}
    \end{equation}
    which is not possible in $\mathbb{Z}/{4\mathbb{Z}}$.\\
    If $\beta$ is odd, then reducing \eqref{eq9} modulo 4 gives,
    \[Y^2 \equiv 3^m \equiv 3 \pmod{4}\] or
    \begin{equation}
        Y^2-Z^2 \equiv 3^m \equiv 3 \pmod{4}
    \end{equation}
    which is not again possible in $\mathbb{Z}/{4\mathbb{Z}}$ and we are done.
    \vspace{2mm}
\end{proof}

\end{document}
